\chapter{Empirische Analyse der Graphklassen unter \textit{stable}-Unabh\"angigkeit}
\label{ch:graphclasses-empirical}

\section{Einleitung}
\label{sec:gce-einleitung}

Kapitel~\ref{ch:graphclasses} hat gezeigt, dass einzelne Graphklassen aus
\emph{berechnungstheoretischer} Sicht sehr unterschiedlich auf die Frage nach
\textit{stable}-Unabh\"angigkeit reagieren: F\"ur azyklische und
geradezykelfreie AFs ist die Antwort strukturell vorbestimmt, w\"ahrend
vollst\"andig gerichtete (und allgemeiner dichte symmetrische) AFs nahezu jede
Anfrage abh\"angig machen. Dieses Kapitel erg\"anzt die theoretische
Einordnung um eine \emph{empirische} Analyse. Ziel ist es, die einzelnen
Klassen anhand konkreter Messdaten genau zu untersuchen und die dabei
wiederkehrenden Muster herauszuarbeiten.

Wir betrachten dazu die Entscheidung von
\(\textit{Ind}_{\mathrm{st}}\) bzw. der Singleton-Variante
\(\textit{SingInd}_{\mathrm{st}}\) (Kapitel~\ref{ch:singl-ci}) auf
Instanzfamilien, die jeweils eine Graphklasse repr\"asentieren. Die zugrunde
liegende Definition bedingter Unabh\"angigkeit ist
Definition~\ref{def:BU:AFs}; wir verwenden durchg\"angig die Konvention, dass
ein positives Ergebnis (SAT) \emph{Unabh\"angigkeit} und ein negatives Ergebnis
(UNSAT) \emph{Abh\"angigkeit} bedeutet.

Die Leitfragen des Kapitels lauten:
\begin{enumerate}
    \item Wie h\"aufig tritt Abh\"angigkeit in den einzelnen Klassen \"uberhaupt
    auf?
    \item Welche strukturellen Gr\"o\ss en --- insbesondere die Anzahl der
    \textit{stable}-Labellings und die Angriffsdichte --- bestimmen diese
    H\"aufigkeit?
    \item L\"asst sich aus den klassenspezifischen Beobachtungen ein
    \"ubergreifendes Muster ableiten, das die theoretischen Resultate aus
    Kapitel~\ref{ch:graphclasses} erkl\"art und vorhersagt?
\end{enumerate}

\section{Versuchsaufbau}
\label{sec:gce-aufbau}

\subsection{Messgr\"o\ss en und Werkzeugkette}
\label{subsec:gce-werkzeug}

Alle Messungen beruhen auf der in Kapitel~\ref{ch:stable} entwickelten
aussagenlogischen Kodierung \(\mathrm{st}_F\) (Definition~\ref{def:st-pl}) der
\textit{stable}-Semantik. Eine Anfrage \((F,X,Y,Z)\) wird in die QBF
\(\Phi^{\mathrm{st}}_F(X,Y\mid Z)\) \"ubersetzt und an den L\"oser QuAbS
\"ubergeben; im Singleton-Fall \(X=\{x\}\), \(Y=\{y\}\) wird zus\"atzlich das
inkrementelle CEGAR-Verfahren aus Kapitel~\ref{ch:singl-ci} verwendet. Zu
jeder Instanz wird die L\"oserantwort (SAT/UNSAT/TIMEOUT) sowie die Laufzeit
protokolliert.

Eine zentrale Begleitgr\"o\ss e ist die \emph{Anzahl der
\textit{stable}-Labellings} \(|\Lambda_{\mathrm{st}}(F)|\). Da diese Zahl
exponentiell gro\ss{} werden kann, erfassen wir sie nur in der Dreiteilung
\[
    n_{\mathrm{st}}(F)\in\{\,0,\ 1,\ {\geq}2\,\},
\]
also \enquote{keine}, \enquote{genau eine} oder \enquote{mindestens zwei}
\textit{stable}-Labellings. Diese Dreiteilung l\"asst sich mit zwei rein
existentiellen Erf\"ullbarkeitsabfragen bestimmen: Eine Abfrage pr\"uft die
Existenz \emph{mindestens einer}, die zweite die Existenz \emph{zweier
verschiedener} \textit{stable}-Labellings. Wie wir in
Abschnitt~\ref{subsec:gce-anker} pr\"azisieren, ist \(n_{\mathrm{st}}(F)\) der
entscheidende Pr\"adiktor f\"ur das Auftreten von Abh\"angigkeit.

\subsection{Messmodi}
\label{subsec:gce-modi}

Wir verwenden zwei komplement\"are Messmodi:

\begin{description}
    \item[Zufallsanfragen.] Pro Instanz wird eine zuf\"allige, reproduzierbar
    geseedete Singleton-Anfrage \(X=\{x\}\), \(Y=\{y\}\) mit \(Z=\emptyset\)
    gestellt und zus\"atzlich \(n_{\mathrm{st}}(F)\) bestimmt. Dieser Modus
    deckt viele Instanzen und mehrere Seeds je Parametrisierung ab und erlaubt
    es, Abh\"angigkeit mit \(n_{\mathrm{st}}\) zu korrelieren.

    \item[Ersch\"opfende Singleton-Enumeration.] F\"ur ausgew\"ahlte
    Repr\"asentanten wird \emph{jedes} ungeordnete Argumentpaar
    \(\{a\},\{b\}\) mit \(a<b\) und \(Z=\emptyset\) getestet. Bei \(n\)
    Argumenten sind das \(\binom{n}{2}\) Anfragen. Dieser Modus liefert ein
    vollst\"andiges Abh\"angigkeitsprofil einer Instanz und ist daher
    unempfindlich gegen\"uber der zuf\"alligen Wahl einzelner Paare.
\end{description}

Der Sonderfall \(Z=\emptyset\) ist f\"ur die Singleton-Enumeration bewusst
gew\"ahlt: Nach den Algorithmen aus Kapitel~\ref{ch:singl-ci} und
Kapitel~\ref{ch:singl-ci-direkt} bricht die Enumeration der gemeinsamen
\(Z\)-Belegungen wegen \(\mathrm{Block}(\eta)=\bot\) nach genau einer
Iteration je Konfiguration ab. Die Einzelabfragen sind dadurch sehr g\"unstig,
was die quadratische Zahl an Paaren \"uberhaupt erst handhabbar macht.

\subsection{Instanzfamilien}
\label{subsec:gce-instanzen}

Tabelle~\ref{tab:gce-familien} fasst die untersuchten Graphklassen und ihre
Erzeugung zusammen. Die symmetrischen AFs werden parametrisch \"uber die
Kantenwahrscheinlichkeit \(p\) erzeugt: Jedes ungeordnete Argumentpaar wird mit
Wahrscheinlichkeit \(p\) durch einen wechselseitigen Angriff verbunden. Der
Grenzfall \(p=1\) liefert genau die vollst\"andig gerichteten AFs aus
Abschnitt~\ref{subsec:cdg-afs}. Die mit \enquote{st} bezeichneten Instanzen
sind durch Teilstichproben (Subsampling) realer Wettbewerbsinstanzen der
ICCMA'25 entstanden und dienen als Stellvertreter f\"ur \enquote{allgemeine},
nicht eigens konstruierte AFs.

\begin{table}[htbp]
    \myfloatalign
    \begin{tabularx}{\textwidth}{lX} \toprule
        Klasse & Erzeugung / Eigenschaft \\ \midrule
        azyklisch            & ohne gerichtete Zyklen; h\"ochstens ein
                               \textit{stable}-Labelling \\
        geradezykelfrei      & ohne Zyklen gerader L\"ange; h\"ochstens ein
                               \textit{stable}-Labelling \\
        ungeradezykelfrei    & ohne Zyklen ungerader L\"ange; stets mindestens
                               ein \textit{stable}-Labelling \\
        SCC-strukturiert     & dichte, stark zusammenh\"angende Komponenten
                               unterschiedlicher Gr\"o\ss e \\
        symmetrisch \((p)\)  & wechselseitige Angriffe mit Wahrscheinlichkeit
                               \(p\in\{0{,}1;0{,}3;0{,}5;0{,}8;0{,}9;1{,}0\}\) \\
        st (ICCMA-Subsample) & Teilstichproben realer Wettbewerbsinstanzen \\
        \bottomrule
    \end{tabularx}
    \caption[Untersuchte Graphklassen]{Untersuchte Graphklassen und ihre
    Erzeugung. Die symmetrischen AFs sind \"uber die Kantenwahrscheinlichkeit
    \(p\) parametrisiert; \(p=1\) entspricht dem vollst\"andig gerichteten AF.}
    \label{tab:gce-familien}
\end{table}

Die Klassen \enquote{geradezykelfrei} und \enquote{ungeradezykelfrei} werden im
Folgenden vor allem theoretisch eingeordnet, da sie sich bereits durch das
Kriterium \(n_{\mathrm{st}}(F)\) vollst\"andig erkl\"aren lassen.

\section{Ein theoretischer Ankerpunkt}
\label{subsec:gce-anker}

Bevor wir die Klassen einzeln betrachten, halten wir ein Resultat fest, das
sich als roter Faden durch alle Messungen zieht. Es verallgemeinert
Proposition~\ref{prop:acyc-af-ci-st-top} von azyklischen AFs auf jede Instanz
mit h\"ochstens einem \textit{stable}-Labelling.

\begin{Proposition}
    \label{prop:gce-leq1-indep}
    Sei \(F=(A,R)\) ein AF mit \(|\Lambda_{\mathrm{st}}(F)|\leq 1\). Dann gilt
    \(X\bot_{\mathrm{st}}Y\mid Z\) f\"ur alle paarweise disjunkten Mengen
    \(X,Y,Z\subseteq A\).
\end{Proposition}

\begin{proof}
    Ist \(\Lambda_{\mathrm{st}}(F)=\emptyset\), so ist die
    \"uber alle Paare quantifizierte Unabh\"angigkeitsbedingung aus
    Definition~\ref{def:BU:AFs} vakuos erf\"ullt. Ist
    \(\Lambda_{\mathrm{st}}(F)=\{\lambda\}\), so gilt f\"ur jedes Paar
    \(\lambda_1,\lambda_2\in\Lambda_{\mathrm{st}}(F)\) bereits
    \(\lambda_1=\lambda_2=\lambda\). Mit \(\lambda_3:=\lambda\) sind
    \(\lambda_3|_X=\lambda_1|_X\), \(\lambda_3|_Y=\lambda_2|_Y\) und
    \(\lambda_3|_Z=\lambda_1|_Z\) trivial erf\"ullt. Also gilt
    \(X\bot_{\mathrm{st}}Y\mid Z\).
\end{proof}

Damit ist \(n_{\mathrm{st}}(F)\geq 2\) eine \emph{notwendige} Bedingung f\"ur
Abh\"angigkeit. Sie ist nicht hinreichend: Auch bei vielen
\textit{stable}-Labellings kann jede konkrete Anfrage unabh\"angig sein
(vgl.\ Korollar~\ref{cor:singl-fixlabel}). Die folgende Analyse zeigt, dass
dieser eine Schwellenwert den gr\"o\ss ten Teil der klassenspezifischen
Unterschiede bereits erkl\"art.

\section{Klassenweise Analyse}
\label{sec:gce-klassen}

\subsection{Azyklische und geradezykelfreie AFs}
\label{subsec:gce-acyc}

F\"ur diese beiden Klassen gilt nach
Proposition~\ref{prop:acyc-af-ci-st-top} bzw.
Proposition~\ref{prop:even-cycle-free-af-ci-st-top} stets
\(|\Lambda_{\mathrm{st}}(F)|\leq 1\). Mit
Proposition~\ref{prop:gce-leq1-indep} ist jede Anfrage unabh\"angig, und zwar
unabh\"angig von der Wahl von \(X\), \(Y\) und \(Z\). Empirisch sind diese
Klassen daher uninteressant im Wortsinn: Es gibt nichts zu finden, weil
Abh\"angigkeit strukturell ausgeschlossen ist. In der Auswertung treten sie
nicht als eigene Messreihe auf, sondern als der theoretisch verstandene
Randfall \(n_{\mathrm{st}}\leq 1\), der in
Abschnitt~\ref{sec:gce-muster} quantitativ best\"atigt wird.

\subsection{Ungeradezykelfreie AFs}
\label{subsec:gce-odd}

Ungeradezykelfreie AFs besitzen stets mindestens ein
\textit{stable}-Labelling, k\"onnen aber auch mehrere besitzen. Anders als bei
den beiden vorigen Klassen gibt es daher keine strukturelle Garantie f\"ur
Unabh\"angigkeit; das Kriterium \(n_{\mathrm{st}}\leq 1\) greift hier im
Allgemeinen nicht. Berechnungstheoretisch bietet die Klasse keinen Vorteil
(Abschnitt~\ref{sec:odd-cycle-free-afs}): Die Bestimmung der
\textit{stable}-Labellings bleibt so schwer wie im allgemeinen Fall, und damit
verbleibt auch \(\textit{Ind}_{\mathrm{st}}\) auf dem vollen
Schwierigkeitsniveau aus Theorem~\ref{th:ind-st-is-pi2p-c}. Die folgenden,
empirisch erschlossenen Klassen lassen sich als konkrete Auspr\"agungen
innerhalb dieses allgemeinen Rahmens lesen.

\subsection{Stark zusammenh\"angende, dichte Komponenten (SCC)}
\label{subsec:gce-scc}

Die SCC-strukturierten Instanzen sind dicht und stark zusammenh\"angend und
umfassen Gr\"o\ss en von einigen hundert bis mehreren tausend Argumenten. Das
empirische Bild ist eindeutig: Die \"uberw\"altigende Mehrheit dieser Instanzen
besitzt \emph{gar kein} \textit{stable}-Labelling.

\begin{table}[htbp]
    \myfloatalign
    \begin{tabularx}{\textwidth}{Xcccc} \toprule
        & \(n_{\mathrm{st}}=0\) & \(n_{\mathrm{st}}=1\) &
          \(n_{\mathrm{st}}\geq2\) & unbestimmt \\ \midrule
        Anzahl Instanzen & 18 & 0 & 1 & 1 \\
        \bottomrule
    \end{tabularx}
    \caption[Stabile Labellings in SCC-Instanzen]{Verteilung von
    \(n_{\mathrm{st}}\) \"uber die 20 ausgewerteten SCC-Instanzen. Die Spalte
    \enquote{unbestimmt} z\"ahlt eine Instanz, deren Z\"ahlabfrage in das
    Zeitlimit lief.}
    \label{tab:gce-scc}
\end{table}

Wie Tabelle~\ref{tab:gce-scc} zeigt, haben 18 der 20 Instanzen kein
\textit{stable}-Labelling; mit Proposition~\ref{prop:gce-leq1-indep} sind sie
damit vakuos unabh\"angig. Entsprechend lieferte \emph{keine} der
ausgewerteten SCC-Anfragen ein abh\"angiges Ergebnis. Dieses Verhalten
best\"atigt die in Abschnitt~\ref{sec:scc-afs} formulierte Beobachtung, dass in
dieser Klasse kaum \textit{stable}-Labellings vorhanden sind. Anschaulich
erzeugen die in dichten Komponenten allgegenw\"artigen ungeraden Angriffszyklen
\enquote{Konflikte ohne Ausweg}: Sie verhindern, dass ein widerspruchsfreies,
alle Argumente entscheidendes Labelling existiert.

Eine praktische Konsequenz betrifft die Laufzeit. Zwar ist jede Antwort
\enquote{unabh\"angig}, doch der Existenznachweis
\(\Lambda_{\mathrm{st}}(F)=\emptyset\) ist auf sehr gro\ss en Instanzen
keineswegs gratis: Die wenigen beobachteten Zeit\"uberschreitungen traten
ausschlie\ss lich bei den gr\"o\ss ten SCC-Instanzen (mehrere tausend
Argumente) auf. Die \emph{Antwort} ist also strukturell vorbestimmt, ihre
\emph{Verifikation} jedoch nicht zwingend billig.

\subsection{Symmetrische AFs: Angriffsdichte als Stellschraube}
\label{subsec:gce-sym}

Die symmetrischen AFs bilden den aufschlussreichsten Teil der Analyse, weil
sich \"uber die Kantenwahrscheinlichkeit \(p\) ein einziger struktureller
Parameter gezielt variieren l\"asst. Nahezu alle erzeugten Instanzen besitzen
mindestens zwei \textit{stable}-Labellings (224 von 225 ausgewerteten
Instanzen), sodass die notwendige Bedingung aus
Proposition~\ref{prop:gce-leq1-indep} praktisch durchg\"angig erf\"ullt ist.
Die Klasse ist damit der ideale Pr\"ufstein f\"ur die Frage, \emph{wie oft}
Abh\"angigkeit tats\"achlich auftritt, sobald sie strukturell m\"oglich ist.

\begin{table}[htbp]
    \myfloatalign
    \begin{tabularx}{\textwidth}{Xcccc} \toprule
        \(p\) & \multicolumn{2}{c}{Singleton-Enumeration} &
                \multicolumn{2}{c}{Zufallsanfragen} \\
        \cmidrule(lr){2-3}\cmidrule(lr){4-5}
              & Paare & abh.\ Anteil & Anfragen & abh.\ Anteil \\ \midrule
        \(0{,}1\) & 4\,950 & 9{,}9\,\%  & 90 & 10{,}0\,\% \\
        \(0{,}3\) & 4\,950 & 28{,}8\,\% & 20 & 35{,}0\,\% \\
        \(0{,}5\) & 1\,225 & 51{,}3\,\% & 44 & 61{,}4\,\% \\
        \(0{,}8\) & 4\,950 & 80{,}7\,\% & 10 & 60{,}0\,\% \\
        \(0{,}9\) & 4\,950 & 90{,}4\,\% & 10 & 90{,}0\,\% \\
        \(1{,}0\) & ---    & ---        & 49 & 100{,}0\,\% \\
        \bottomrule
    \end{tabularx}
    \caption[Abh\"angigkeitsrate symmetrischer AFs nach Dichte]{Anteil
    abh\"angiger Anfragen in symmetrischen AFs in Abh\"angigkeit von der
    Kantenwahrscheinlichkeit \(p\). Die Singleton-Enumeration testet
    ersch\"opfend alle \(\binom{n}{2}\) Paare eines Repr\"asentanten (\(n=100\),
    f\"ur \(p=0{,}5\) ein Repr\"asentant mit \(n=50\)); die Zufallsanfragen
    aggregieren mehrere Instanzen und Seeds.}
    \label{tab:gce-sym}
\end{table}

Tabelle~\ref{tab:gce-sym} zeigt das zentrale Ergebnis: Der Anteil abh\"angiger
Anfragen w\"achst \emph{monoton} mit der Dichte \(p\). In der ersch\"opfenden
Enumeration steigt er von \(9{,}9\,\%\) bei \(p=0{,}1\) \"uber rund die H\"alfte
bei \(p=0{,}5\) auf \(90{,}4\,\%\) bei \(p=0{,}9\). Die unabh\"angig erhobenen
Zufallsanfragen best\"atigen denselben Verlauf und kulminieren bei \(p=1\) in
einer Abh\"angigkeitsrate von \(100\,\%\). Letzteres ist kein Zufall, sondern
exakt die Aussage von Proposition~\ref{prop:cdg-afs-not-st-ci}: Im
vollst\"andig gerichteten AF kann jedes Argument nur \emph{allein}
akzeptiert werden, weshalb sich die Akzeptanz zweier verschiedener Argumente
niemals kombinieren l\"asst. Die symmetrische Familie interpoliert somit
stetig zwischen dem \enquote{leichten} Randfall d\"unner Graphen und dem
vollst\"andig abh\"angigen Extremfall.

Die strukturelle Erkl\"arung liefert Abschnitt~\ref{sec:sym-afs}: Jeder
zus\"atzliche (wechselseitige) Angriff schr\"ankt die Kombinierbarkeit von
Labellings ein. Mehr Kanten bedeuten mehr Paare \((\lambda_1,\lambda_2)\),
deren \enquote{Kreuzkombination} \(\lambda_3\) durch einen Konflikt blockiert
wird --- und genau ein solches Paar ist nach
Kapitel~\ref{ch:singl-ci-direkt} ein Gegenbeispiel zur Unabh\"angigkeit.

\subsection{Allgemeine AFs (ICCMA-Subsamples)}
\label{subsec:gce-st}

Die st-Instanzen repr\"asentieren \enquote{nat\"urlich} vorkommende AFs, die
nicht auf eine bestimmte zyklische Eigenschaft hin konstruiert sind. Ihr
Verhalten liegt erwartungsgem\"a\ss{} zwischen den Extremen.

\begin{table}[htbp]
    \myfloatalign
    \begin{tabularx}{\textwidth}{Xcccc} \toprule
        & \(n_{\mathrm{st}}=0\) & \(n_{\mathrm{st}}=1\) &
          \(n_{\mathrm{st}}\geq2\) & unbestimmt \\ \midrule
        Anzahl Anfragen & 9 & 24 & 71 & 4 \\
        \bottomrule
    \end{tabularx}
    \caption[Stabile Labellings in st-Instanzen]{Verteilung von
    \(n_{\mathrm{st}}\) \"uber die ausgewerteten st-Zufallsanfragen. Ein Drittel
    der Instanzen besitzt h\"ochstens ein \textit{stable}-Labelling und ist
    damit nach Proposition~\ref{prop:gce-leq1-indep} garantiert unabh\"angig.}
    \label{tab:gce-st}
\end{table}

Tabelle~\ref{tab:gce-st} zeigt eine breite Streuung von \(n_{\mathrm{st}}\):
33 von 108 Anfragen entfallen auf Instanzen mit \(n_{\mathrm{st}}\leq 1\), die
nach Proposition~\ref{prop:gce-leq1-indep} unabh\"angig sind. Insgesamt waren
die st-Zufallsanfragen ganz \"uberwiegend unabh\"angig (102 SAT gegen\"uber nur
2 UNSAT, bei 4 Zeit\"uberschreitungen); beide abh\"angigen F\"alle traten auf
Instanzen mit \(n_{\mathrm{st}}\geq2\) auf.

Das vollst\"andige Abh\"angigkeitsprofil zweier Repr\"asentanten unterstreicht,
dass Abh\"angigkeit in dieser Klasse zwar selten, aber real ist. F\"ur die
Instanz \texttt{st\_156} sind von \(12\,090\) getesteten Paaren \(4{,}5\,\%\)
abh\"angig; f\"ur \texttt{st\_211} sind es von \(22\,155\) Paaren
\(10{,}6\,\%\). Reale AFs verhalten sich also eher wie d\"unne symmetrische
Graphen (kleiner \(p\)) als wie dichte: Die meisten Argumentpaare sind
unabh\"angig, einzelne lokale Konfliktstrukturen erzeugen jedoch gezielte
Abh\"angigkeiten.

Bemerkenswert ist schlie\ss lich die Laufzeitspreizung innerhalb der Klasse.
W\"ahrend die meisten Anfragen praktisch sofort entschieden werden, traten
sowohl die Zeit\"uberschreitungen als auch deutliche Ausrei\ss er (eine Instanz
mit Laufzeiten im Bereich von rund zwei bis vier Minuten) gerade hier auf. Die
\(\Pi_2^{\mathrm P}\)-H\"arte aus Theorem~\ref{th:ind-st-is-pi2p-c} wird damit
nicht von den konstruierten Extremklassen, sondern von den
\enquote{allgemeinen} Instanzen sichtbar gemacht.

\section{\"Ubergreifende Muster}
\label{sec:gce-muster}

Aus der klassenweisen Analyse lassen sich vier wiederkehrende Muster
verdichten.

\paragraph{Muster 1: Abh\"angigkeit erfordert mindestens zwei
\textit{stable}-Labellings.}
\"Uber alle Klassen hinweg ist \(n_{\mathrm{st}}(F)\geq 2\) ausnahmslos
notwendig f\"ur ein abh\"angiges Ergebnis. Tabelle~\ref{tab:gce-master}
schl\"usselt die Zufallsanfragen nach \(n_{\mathrm{st}}\) auf: In keiner
einzigen Anfrage mit \(n_{\mathrm{st}}\leq 1\) trat Abh\"angigkeit auf, was
Proposition~\ref{prop:gce-leq1-indep} exakt widerspiegelt. Erst im Regime
\(n_{\mathrm{st}}\geq2\) wird Abh\"angigkeit \"uberhaupt m\"oglich --- und dann
auch h\"aufig.

\begin{table}[htbp]
    \myfloatalign
    \begin{tabularx}{\textwidth}{Xccc} \toprule
        \(n_{\mathrm{st}}(F)\) & SAT (unabh.) & UNSAT (abh.) & TIMEOUT \\
        \midrule
        \(0\)       &  28 &   0 & 0 \\
        \(1\)       &  24 &   0 & 0 \\
        \({\geq}2\) & 185 & 110 & 2 \\
        \bottomrule
    \end{tabularx}
    \caption[Abh\"angigkeit nach Anzahl stabiler Labellings]{Zufallsanfragen
    aufgeschl\"usselt nach \(n_{\mathrm{st}}(F)\). Abh\"angigkeit (UNSAT) tritt
    ausschlie\ss lich bei \(n_{\mathrm{st}}\geq2\) auf.}
    \label{tab:gce-master}
\end{table}

\paragraph{Muster 2: Angriffsdichte steuert die H\"aufigkeit.}
Sobald gen\"ugend \textit{stable}-Labellings vorhanden sind, bestimmt die
Angriffsdichte, \emph{wie oft} Abh\"angigkeit auftritt. Die symmetrische
Familie macht dies quantitativ greifbar (Tabelle~\ref{tab:gce-sym}): Die
Abh\"angigkeitsrate w\"achst monoton mit \(p\) und erreicht im dichtesten Fall
\(100\,\%\). Dichte und Abh\"angigkeit sind damit gleichgerichtet --- ein
empirisches Gegenst\"uck zur theoretischen Beobachtung aus
Abschnitt~\ref{sec:sym-afs}, dass zus\"atzliche Angriffe die Kombinierbarkeit
von Labellings einschr\"anken.

\paragraph{Muster 3: Die Klassen ordnen sich in einem Spektrum an.}
Die untersuchten Klassen bilden ein durchgehendes Spektrum von
\enquote{strukturell unabh\"angig} bis \enquote{strukturell abh\"angig}:
\begin{center}
\begin{tabularx}{\textwidth}{lX} \toprule
    unabh\"angiger Pol & azyklisch und geradezyklenfrei
        (\(n_{\mathrm{st}}\leq1\)); dichte SCC
        (\(\Lambda_{\mathrm{st}}=\emptyset\)) \\ \midrule
    Mitte & allgemeine (st) und d\"unne symmetrische AFs \\ \midrule
    abh\"angiger Pol & dichte symmetrische bis vollst\"andig gerichtete AFs
        (\(p\to1\)) \\ \bottomrule
\end{tabularx}
\end{center}
Am einen Ende erzwingen \(n_{\mathrm{st}}\leq 1\) (azyklisch, geradezyklenfrei)
bzw.\ \(\Lambda_{\mathrm{st}}=\emptyset\) (dichte SCC) generische
Unabh\"angigkeit. Am anderen Ende erzwingt der vollst\"andig gerichtete Graph
generische Abh\"angigkeit. Reale (st) und d\"unne symmetrische Instanzen liegen
nahe beim unabh\"angigen Ende, dichte symmetrische Instanzen nahe beim
abh\"angigen.

\paragraph{Muster 4: Antwort und Aufwand sind zu entkoppeln.}
Eine strukturell vorbestimmte Antwort bedeutet nicht automatisch eine billige
Berechnung. Die g\"unstigsten Messungen waren die Singleton-Tests mit
\(Z=\emptyset\), die nach einer Iteration abbrechen
(Abschnitt~\ref{subsec:gce-modi}). Die teuersten F\"alle --- Zeit\"uberschreitungen
und Minuten-Ausrei\ss er --- traten dort auf, wo entweder die Instanzen sehr
gro\ss{} (SCC) oder strukturell unregelm\"a\ss ig (st) waren. Das praktische
Schwierigkeitsbild folgt damit nicht der Abh\"angigkeitsrate, sondern der
Gr\"o\ss e und Regelm\"a\ss igkeit der Instanz.

\section{Zusammenfassung}
\label{sec:gce-zusammenfassung}

Die empirische Analyse best\"atigt und pr\"azisiert die theoretische
Einordnung aus Kapitel~\ref{ch:graphclasses}. Der einzelne Pr\"adiktor
\(n_{\mathrm{st}}(F)\) trennt die Klassen sauber in ein Regime garantierter
Unabh\"angigkeit (\(n_{\mathrm{st}}\leq 1\): azyklisch, geradezyklenfrei, fast
alle SCC-Instanzen) und ein Regime m\"oglicher Abh\"angigkeit
(\(n_{\mathrm{st}}\geq 2\): symmetrisch, allgemein). Innerhalb des zweiten
Regimes ist die Angriffsdichte die entscheidende Stellschraube: Die
symmetrische Familie zeigt einen monotonen Anstieg der Abh\"angigkeitsrate von
rund einem Zehntel bei \(p=0{,}1\) bis zur vollst\"andigen Abh\"angigkeit bei
\(p=1\). Allgemeine ICCMA-Instanzen verhalten sich \"ahnlich wie d\"unne
symmetrische Graphen: \"uberwiegend unabh\"angig, mit vereinzelten, lokal
erzeugten Abh\"angigkeiten. Schlie\ss lich entkoppeln sich strukturelle Antwort
und Rechenaufwand --- die \(\Pi_2^{\mathrm P}\)-H\"arte des Problems wird gerade
auf den allgemeinen, nicht eigens konstruierten Instanzen sichtbar.
